% Shared main-text method for both active manuscripts.
\section{Preliminaries and Self-Harness Optimization}
\label{sec:method}

An \rlm{} drives a persistent Python REPL to inspect an input, issue model calls,
and assemble an answer~\citep{zhang2025rlm}. The input is initially stored in
\texttt{context}; code can expose selected content through printed output. The
ordinary submission path sets \texttt{answer["content"]} and
\texttt{answer["ready"]}. Model weights, experiment-owned resource caps, and
the environment verifier remain fixed during optimization.

Following \citet{zhang2026selfharness}, we expose ten editable harness surfaces:
five instruction fields, runtime policy, execution-output formatting, REPL
helpers, answer middleware, and a skill library. Each surface corresponds to a
declared field or group of fields. Human-authored capability descriptions and
mechanism-to-surface routes constrain the search. The current study starts from
\texttt{H0*R}, which already contains hand-authored recursion and decomposition
guidance; the sparse \hzero{} is a reference baseline. Appendix~\ref{app:surfaces}
specifies the actual inputs, scope, and invocation conditions of each surface.

\subsection{Self-Harness Optimization}
\label{sec:sho}

The same fixed model executes tasks, diagnoses failures, and proposes edits.
Each round has three stages (Figure~\ref{fig:sh-loop}).

\paragraph{Weakness mining.}
The incumbent runs on mining tasks. The environment verifier evaluates final
answers, and an attributor uses selected execution-trace operations to infer
failure mechanisms with explicit uncertainty. Diagnoses use a predefined
vocabulary and are clustered by their recorded failure signature. Cited
operations support an interpretation; they do not independently establish its
causal correctness.

\paragraph{Harness proposal.}
The proposer receives compact operation-level evidence, the incumbent surfaces,
runtime capabilities, and bounded history of attempted edits and outcomes. It
selects an intervention before writing its replacement text and must explain
what the incumbent does, what observed behavior remains wrong, and what changes
at the relevant operation. At most one admitted edit may occupy a surface or
target a pattern in a round. The edit ceiling permits fewer edits or abstention.
Host materialization and preflight reject malformed or structurally invalid
changes before task evaluation.

\paragraph{Joint validation and promotion.}
All admitted edits are composed into one candidate, which is compared with a
freshly evaluated incumbent on the same validation tasks. Validation tasks are
excluded from mining but reused for adaptive selection; aggregate outcomes
inform later proposal history. In the current configuration, promotion requires
at least as many environment-verifier passes as the incumbent and compliance
with the configured resource bands. Optional verifier-defined dense metrics
inform diagnosis and history, but are not promotion conditions. The combined
candidate is the measured unit; its outcome does not identify each edit's
individual contribution. Optimization stops at the round ceiling or after the
configured number of rounds without promotion, then freezes the final incumbent
for separate evaluation. Appendices~\ref{app:mining}--\ref{app:measurement}
give the implementation contracts and measurement definitions.
